\heading{热学-24春-期末2}
\noteinfo{题目来源于赛艇，答案由AI生成。}

\begin{question}
写出热力学第一定律，并说明它与能量守恒的关系。
\begin{solution}
第一定律可写为
\[
dU=\delta Q-\delta W.
\]
对有限过程，
\[
\Delta U=Q-W.
\]
含义是：系统内能的增量等于外界传给系统的热量减去系统对外所作的功。它是能量守恒定律在热学中的具体表达，其中热和功是能量传递的两种方式；但第一定律本身不说明过程方向，这要靠第二定律。
\ansmath{\Delta U=Q-W}
\end{solution}
\end{question}

\begin{question}
写出熵表述的第二定律，并由此证明卡诺定理。
\begin{solution}
熵表述为：任意循环满足
\[
\oint \frac{\delta Q}{T}\le 0,
\]
等号仅对可逆循环成立。设热机从高温热源 $T_h$ 吸热 $Q_h$，向低温热源 $T_c$ 放热 $Q_c$，则
\[
\frac{Q_h}{T_h}-\frac{Q_c}{T_c}\le 0
\quad\Rightarrow\quad
\frac{Q_c}{Q_h}\ge \frac{T_c}{T_h}.
\]
因此效率
\[
\eta=1-\frac{Q_c}{Q_h}\le 1-\frac{T_c}{T_h}.
\]
等号仅对可逆机成立，所以在同样两热源之间，没有任何热机效率能超过可逆卡诺机；所有可逆机效率相同，这就是卡诺定理。
\ansmath{\eta=1-\frac{Q_c}{Q_h}\le 1-\frac{T_c}{T_h}}
\end{solution}
\end{question}

\begin{question}
硬币插入干冰会“发抖”，钢球压在干冰上会“尖叫”，为什么？
\begin{solution}
干冰会迅速升华。硬币或钢球与干冰接触后，其下方不断产生二氧化碳气体，形成很薄的气膜。气膜的喷射与局部接触会引起周期性的滑移/振动：
\begin{itemize}[leftmargin=1.8em,itemsep=0.2em]
\item 硬币较轻，振动频率较低，表现为“发抖”；
\item 钢球压得更紧，气体从狭缝高速喷出，更容易激发高频机械振动，故会发出尖锐声音。
\end{itemize}
本质上都是升华气体导致的自激振动与声辐射。
\anstext{见上。}
\end{solution}
\end{question}

\begin{question}
黑色陶瓷和白色陶瓷一起加热到约 $1000^\circ\mathrm C$ 时看起来分别是什么颜色？可总结什么规律？
\begin{solution}
两者都会发出明显热辐射，肉眼看上去都接近暗红到红橙色；黑色陶瓷通常更亮一些，但颜色主调主要由温度决定，而不是室温下的反射颜色。可总结：\textbf{热辐射光谱主要由温度决定；吸收率高的材料往往发射率也高}（Kirchhoff 定律）。
\anstext{两者都会发出明显热辐射，肉眼看上去都接近暗红到红橙色；黑色陶瓷通常更亮一些，但颜色主调主要由温度决定，而不是室温下的反射颜色。可总结：（Kirchhoff 定律）。}
\end{solution}
\end{question}

\begin{question}
理想气体 Brayton 循环：$1\to2,3\to4$ 为绝热，$2\to3,4\to1$ 为等压。请提出至少两种提高热效率的方法。
\begin{solution}
对理想 Brayton 循环，热效率可写成
\[
\eta=1-\frac{T_4-T_1}{T_3-T_2}.
\]
提高效率的常见办法有：
\begin{itemize}[leftmargin=1.8em,itemsep=0.2em]
\item 提高压比 $p_2/p_1$；
\item 提高涡轮进口温度 $T_3$；
\item 加装回热器（再生器），用尾气预热压缩后的工质；
\item 采用多级压缩配中间冷却、或多级膨胀配再热并优化参数。
\end{itemize}
只要答出其中两条并说明其物理原因即可。
\ansmath{\eta=1-\frac{T_4-T_1}{T_3-T_2}}
\end{solution}
\end{question}

\begin{question}
在绝热容器中，质量都为 $m$ 的两份水，初温分别为 $T_1,T_2$，压强一定，比热 $C_p$ 取常数。
\begin{enumerate}[label=（\arabic*）,leftmargin=2.8em,itemsep=0.2em]
\item 求总熵变；
\item 对每一份水分别看，该过程是否可逆？第二定律怎样表述才正确？
\item 若 $T_1>T_2$，利用混合前两份水作为热源，最多可输出多少功？
\end{enumerate}
\begin{solution}
设直接混合后的平衡温度为
\[
T_f=\frac{T_1+T_2}{2}.
\]
\begin{enumerate}[label=（\arabic*）,leftmargin=2.8em,itemsep=0.2em]
\item 总熵变
\[
\Delta S=mC_p\ln\frac{T_f}{T_1}+mC_p\ln\frac{T_f}{T_2}
=mC_p\ln\frac{(T_1+T_2)^2}{4T_1T_2}\ge 0.
\]
\item 对每一部分水而言，实际换热都跨越有限温差，因此都是\textbf{不可逆}过程。若把单独一部分水拿出来看，它不是孤立系，故不能直接说“熵不减”；正确说法是 Clausius 不等式适用于实际过程，而对整个绝热复合系统则有总熵增加。
\item 若通过可逆热机榨取最大功，则总熵守恒，末态共同温度满足
\[
2mC_p\ln T_f^{(\max W)}=mC_p\ln T_1+mC_p\ln T_2,
\]
故
\[
T_f^{(\max W)}=\sqrt{T_1T_2}.
\]
最大输出功为
\[
W_{\max}=mC_p(T_1+T_2-2\sqrt{T_1T_2})
=mC_p(\sqrt{T_1}-\sqrt{T_2})^2.
\]
\end{enumerate}
\anstext{见上。}
\end{solution}
\end{question}

\begin{question}
水的气液相变。
\begin{enumerate}[label=（\arabic*）,leftmargin=2.8em,itemsep=0.2em]
\item 用熵判据推导气液共存时的平衡条件；
\item 饱和蒸气视为理想气体，且 $v_g\gg v_l$、汽化热 $L$ 为常量，求饱和蒸气压方程。
\end{enumerate}
\begin{solution}
\begin{enumerate}[label=（\arabic*）,leftmargin=2.8em,itemsep=0.2em]
\item 对孤立两相系统，平衡时总熵极大。对两相间微小交换 $dU,dV,dN$，有
\[
dS_{\text{tot}}=\left(\frac1{T_g}-\frac1{T_l}\right)dU+
\left(\frac{p_g}{T_g}-\frac{p_l}{T_l}\right)dV-
\left(\frac{\mu_g}{T_g}-\frac{\mu_l}{T_l}\right)dN.
\]
平衡时对任意允许扰动都应为零，故
\[
T_g=T_l,\qquad p_g=p_l,\qquad \mu_g=\mu_l.
\]
\item Clausius--Clapeyron 方程为
\[
\frac{dp}{dT}=\frac{L}{T(v_g-v_l)}\approx \frac{L}{T v_g}.
\]
又理想气体近似 $v_g=RT/p$，故
\[
\frac{dp}{dT}=\frac{Lp}{RT^2}.
\]
分离变量积分得
\[
\ln p=-\frac{L}{RT}+C,
\]
即
\[
p(T)=p_*\exp\left[-\frac{L}{R}\left(\frac1T-\frac1{T_*}\right)\right],
\]
或写作 $p=Ae^{-L/(RT)}$。
\end{enumerate}
\anstext{见上。}
\end{solution}
\end{question}

\begin{question}
光子气体：已知压强 $p=u/3$，且内能密度 $u$ 只依赖于温度 $T$。
\begin{enumerate}[label=（\arabic*）,leftmargin=2.8em,itemsep=0.2em]
\item 推导 $U$ 与 $T,V$ 的关系；
\item 推导绝热过程中 $V,T$ 的关系。
\end{enumerate}
\begin{solution}
设总内能
\[
U=Vu(T).
\]
由
\[
dS=\frac{1}{T}dU+\frac{p}{T}dV
=\frac{V u'(T)}{T}dT+\frac{u+p}{T}dV
=\frac{V u'(T)}{T}dT+\frac{4u}{3T}dV.
\]
由 $S(T,V)$ 的全微分可积条件，得
\[
\frac{u'}{T}=\frac{d}{dT}\left(\frac{4u}{3T}\right)
\quad\Rightarrow\quad
u'(T)=\frac{4u(T)}{T}.
\]
解得
\[
u(T)=aT^4,
\qquad U=aVT^4.
\]
绝热可逆时 $dS=0$，等价于
\[
dU+p\,dV=0.
\]
代入 $U=aVT^4$ 与 $p=aT^4/3$，得
\[
4aVT^3dT+\frac{4}{3}aT^4dV=0
\quad\Rightarrow\quad
\frac{dT}{T}=-\frac13\frac{dV}{V}.
\]
故
\[
TV^{1/3}=\text{常数}.
\]
\ansmath{TV^{1/3}=\text{常数}}
\end{solution}
\end{question}

\begin{question}
半径 $r=1\,\mathrm{mm}$ 的气泡位于水下 $h=10\,\mathrm m$，水温 $T_0=300\,\mathrm K$，表面张力系数 $\alpha=0.073\,\mathrm{N/m}$，泡内初压约为 $10^{-5}$ 个大气压，估算内爆时最高温度；可能出现什么现象？设水蒸气 $C_V=3R$。
\begin{solution}
取外压为
\[
p_f\approx p_0+\rho gh+\frac{2\alpha}{r}
\approx 1\,\text{atm}+0.97\,\text{atm}+0.0014\,\text{atm}
\approx 2.0\,\text{atm}.
\]
初压
\[
p_i\approx 10^{-5}\,\text{atm}.
\]
因为 $C_V=3R$，故
\[
\gamma=\frac{C_p}{C_V}=\frac{4R}{3R}=\frac43.
\]
把内爆近似为绝热压缩，
\[
T_f=T_0\left(\frac{p_f}{p_i}\right)^{(\gamma-1)/\gamma}
=300\left(\frac{2.0}{10^{-5}}\right)^{1/4}
\approx 6.3\times 10^3\,\mathrm K.
\]
量级约为数千开。如此高温下可能出现强烈发光（声致发光）、局部电离等现象。
\ansmath{T_f=T_0\left(\frac{p_f}{p_i}\right)^{(\gamma-1)/\gamma}
=300\left(\frac{2.0}{10^{-5}}\right)^{1/4}
\approx 6.3\times 10^3\,\mathrm K}
\end{solution}
\end{question}


